% !TEX root = ../main.tex
% =====================================================================
%  CHƯƠNG 5 — THIẾT KẾ VÀ THI CÔNG PHẦN CỨNG
% ---------------------------------------------------------------------
%  Người phụ trách : [TODO: Họ tên - MSSV]
%  Tiền tố label   : c5      (fig:c5:..., tab:c5:..., eq:c5:..., sec:c5:...)
%  Thư mục ảnh     : image/chuong5/
%  File TLTK       : references/chuong5.bib   (key: c5_...)
%  Độ dài gợi ý    : 5–10 trang
%  Mục tiêu        : Schematic (nguồn + chức năng), layout PCB, ảnh mạch thực
%                    tế, các lỗi gặp phải & cách khắc phục, đánh giá phần cứng.
% =====================================================================
\chapter{Thiết kế và thi công phần cứng}
\label{chap:c5}

% ---------------------------------------------------------------------
\section{Sơ đồ nguyên lý khối nguồn}
\label{sec:c5:khoi-nguon}

\TODO{Trình bày các mức nguồn cần thiết, IC nguồn được chọn, schematic.}

% ---------------------------------------------------------------------
\section{Sơ đồ nguyên lý khối chức năng}
\label{sec:c5:khoi-chuc-nang}

\TODO{Schematic từng khối: khuếch đại, bù nhiệt, V-I, đọc dòng, ADC, connector.}

% ---------------------------------------------------------------------
\section{Thiết kế mạch in (PCB)}
\label{sec:c5:layout}
% Gợi ý: số lớp, độ rộng đường mạch, cách đi dây analog/digital, mặt top/bottom, 3D.

\TODO{Trình bày layout PCB, ảnh top/bottom/3D.}

% ---------------------------------------------------------------------
\section{Phần cứng thực tế}
\label{sec:c5:thuc-te}

\TODO{Ảnh mạch sau khi hàn, các lỗi phát sinh và cách khắc phục.}

% ---------------------------------------------------------------------
\section{Đánh giá phần cứng}
\label{sec:c5:danh-gia}

\subsection{Ưu điểm}
\label{sec:c5:uu-diem}
\begin{itemize}
    \item \TODO{}
\end{itemize}

\subsection{Hạn chế}
\label{sec:c5:han-che}
\begin{itemize}
    \item \TODO{}
\end{itemize}
